% Options for packages loaded elsewhere
% Options for packages loaded elsewhere
\PassOptionsToPackage{unicode}{hyperref}
\PassOptionsToPackage{hyphens}{url}
%
\documentclass[
  ignorenonframetext,
  aspectratio=32,
]{beamer}
\newif\ifbibliography
\usepackage{pgfpages}
\setbeamertemplate{caption}[numbered]
\setbeamertemplate{caption label separator}{: }
\setbeamercolor{caption name}{fg=normal text.fg}
\beamertemplatenavigationsymbolshorizontal
% remove section numbering
\setbeamertemplate{part page}{
  \centering
  \begin{beamercolorbox}[sep=16pt,center]{part title}
    \usebeamerfont{part title}\insertpart\par
  \end{beamercolorbox}
}
\setbeamertemplate{section page}{
  \centering
  \begin{beamercolorbox}[sep=12pt,center]{section title}
    \usebeamerfont{section title}\insertsection\par
  \end{beamercolorbox}
}
\setbeamertemplate{subsection page}{
  \centering
  \begin{beamercolorbox}[sep=8pt,center]{subsection title}
    \usebeamerfont{subsection title}\insertsubsection\par
  \end{beamercolorbox}
}
% Prevent slide breaks in the middle of a paragraph
\widowpenalties 1 10000
\raggedbottom
\AtBeginPart{
  \frame{\partpage}
}
\AtBeginSection{
  \ifbibliography
  \else
    \frame{\sectionpage}
  \fi
}
\AtBeginSubsection{
  \frame{\subsectionpage}
}
\usepackage{iftex}
\ifPDFTeX
  \usepackage[T1]{fontenc}
  \usepackage[utf8]{inputenc}
  \usepackage{textcomp} % provide euro and other symbols
\else % if luatex or xetex
  \usepackage{unicode-math} % this also loads fontspec
  \defaultfontfeatures{Scale=MatchLowercase}
  \defaultfontfeatures[\rmfamily]{Ligatures=TeX,Scale=1}
\fi
\usepackage{lmodern}

\usetheme[]{Copenhagen}
\usecolortheme{lily}
\ifPDFTeX\else
  % xetex/luatex font selection
\fi
% Use upquote if available, for straight quotes in verbatim environments
\IfFileExists{upquote.sty}{\usepackage{upquote}}{}
\IfFileExists{microtype.sty}{% use microtype if available
  \usepackage[]{microtype}
  \UseMicrotypeSet[protrusion]{basicmath} % disable protrusion for tt fonts
}{}
\makeatletter
\@ifundefined{KOMAClassName}{% if non-KOMA class
  \IfFileExists{parskip.sty}{%
    \usepackage{parskip}
  }{% else
    \setlength{\parindent}{0pt}
    \setlength{\parskip}{6pt plus 2pt minus 1pt}}
}{% if KOMA class
  \KOMAoptions{parskip=half}}
\makeatother


\usepackage{longtable,booktabs,array}
\usepackage{calc} % for calculating minipage widths
\usepackage{caption}
% Make caption package work with longtable
\makeatletter
\def\fnum@table{\tablename~\thetable}
\makeatother
\usepackage{graphicx}
\makeatletter
\newsavebox\pandoc@box
\newcommand*\pandocbounded[1]{% scales image to fit in text height/width
  \sbox\pandoc@box{#1}%
  \Gscale@div\@tempa{\textheight}{\dimexpr\ht\pandoc@box+\dp\pandoc@box\relax}%
  \Gscale@div\@tempb{\linewidth}{\wd\pandoc@box}%
  \ifdim\@tempb\p@<\@tempa\p@\let\@tempa\@tempb\fi% select the smaller of both
  \ifdim\@tempa\p@<\p@\scalebox{\@tempa}{\usebox\pandoc@box}%
  \else\usebox{\pandoc@box}%
  \fi%
}
% Set default figure placement to htbp
\def\fps@figure{htbp}
\makeatother





\setlength{\emergencystretch}{3em} % prevent overfull lines

\providecommand{\tightlist}{%
  \setlength{\itemsep}{0pt}\setlength{\parskip}{0pt}}



 


\makeatletter
\@ifpackageloaded{caption}{}{\usepackage{caption}}
\AtBeginDocument{%
\ifdefined\contentsname
  \renewcommand*\contentsname{Table of contents}
\else
  \newcommand\contentsname{Table of contents}
\fi
\ifdefined\listfigurename
  \renewcommand*\listfigurename{List of Figures}
\else
  \newcommand\listfigurename{List of Figures}
\fi
\ifdefined\listtablename
  \renewcommand*\listtablename{List of Tables}
\else
  \newcommand\listtablename{List of Tables}
\fi
\ifdefined\figurename
  \renewcommand*\figurename{Figure}
\else
  \newcommand\figurename{Figure}
\fi
\ifdefined\tablename
  \renewcommand*\tablename{Table}
\else
  \newcommand\tablename{Table}
\fi
}
\@ifpackageloaded{float}{}{\usepackage{float}}
\floatstyle{ruled}
\@ifundefined{c@chapter}{\newfloat{codelisting}{h}{lop}}{\newfloat{codelisting}{h}{lop}[chapter]}
\floatname{codelisting}{Listing}
\newcommand*\listoflistings{\listof{codelisting}{List of Listings}}
\makeatother
\makeatletter
\makeatother
\makeatletter
\@ifpackageloaded{caption}{}{\usepackage{caption}}
\@ifpackageloaded{subcaption}{}{\usepackage{subcaption}}
\makeatother

\usepackage{bookmark}
\IfFileExists{xurl.sty}{\usepackage{xurl}}{} % add URL line breaks if available
\urlstyle{same}
\hypersetup{
  pdftitle={Diurnal Temperature Range Trends in St.~John's, NL (2000--2025)},
  pdfauthor={Isaac Adoboe},
  hidelinks,
  pdfcreator={LaTeX via pandoc}}



\title{Diurnal Temperature Range Trends in St.~John's, NL (2000--2025)}
\author{Isaac Adoboe}
\date{}

\begin{document}
\frame{\titlepage}


\section{Motivation \& Context}\label{motivation-context}

\begin{frame}{Motivation \& Context}
\begin{itemize}
\tightlist
\item
  \textbf{Diurnal Temperature Range (DTR)} =
  \(T_{\max}(t) - T_{\min}(t)\)
\item
  DTR reflects atmospheric stability, cloud cover, and air mass
  characteristics
\item
  \textbf{Global pattern:} DTR declining due to asymmetric warming
\item
  \textbf{St.~John's context:} Maritime climate; winter DTR suppressed,
  summer DTR elevated
\item
  \textbf{Research question:} Is DTR declining in St.~John's consistent
  with global warming theory?
\end{itemize}
\end{frame}

\section{Data \& Methods}\label{data-methods}

\begin{frame}{Data \& Methods}
\textbf{Data Sources} - Four meteorological stations in St.~John's
region - Daily minimum and maximum temperatures (2000--2025) - Regional
mean DTR computed across stations

\textbf{Statistical Approach} - Linear regression:
\(\text{DTR}_{\text{annual}} \sim \text{Year}\) - Mann-Kendall test:
non-parametric test for monotonic trend - Extreme events: High-DTR
(\textgreater90th percentile), Low-DTR (\textless10th percentile) -
Sensitivity: compare global vs.~seasonal threshold definitions -
Significance level: \(\alpha = 0.05\)
\end{frame}

\section{Exploratory Analysis: Seasonal DTR
Pattern}\label{exploratory-analysis-seasonal-dtr-pattern}

\begin{frame}{Exploratory Analysis: Seasonal DTR Pattern}
\begin{figure}[H]

{\centering \includegraphics[width=0.8\linewidth,height=\textheight,keepaspectratio]{../plots/dtr_by_month.pdf}

}

\caption{Monthly mean DTR (placeholder)}

\end{figure}%

\textbf{Key observation:} - Winter DTR suppressed
(\textasciitilde8--10°C) due to maritime moderation - Summer DTR
elevated (\textasciitilde14--16°C) under clearer skies - Pattern
consistent with seasonal cloud cover and radiation effects
\end{frame}

\section{Annual Mean DTR Trend}\label{annual-mean-dtr-trend}

\begin{frame}{Annual Mean DTR Trend}
\begin{figure}[H]

{\centering \includegraphics[width=0.85\linewidth,height=\textheight,keepaspectratio]{../plots/mann_kendall_annual_dtr.pdf}

}

\caption{Mann-Kendall test for annual mean DTR}

\end{figure}%

\textbf{Result:} - Kendall's \(\tau\) = {[}YOUR VALUE{]} - p-value
(one-sided) = {[}YOUR VALUE{]} - \textbf{Conclusion:} {[}No significant
trend{]} / {[}Significant declining trend{]} / {[}Significant increasing
trend{]} - \textbf{Implication:} {[}Consistent with / Contradicts{]}
global warming hypothesis
\end{frame}

\section{Extreme Event Trends}\label{extreme-event-trends}

\begin{frame}{Extreme Event Trends}
\begin{figure}[H]

{\centering \includegraphics[width=0.9\linewidth,height=\textheight,keepaspectratio]{../plots/mann_kendall_extremes.pdf}

}

\caption{High and low DTR extremes over time}

\end{figure}%

\begin{longtable}[]{@{}llll@{}}
\toprule\noalign{}
Extreme Type & Kendall \(\tau\) & p-value & Significant? \\
\midrule\noalign{}
\endhead
High-DTR extremes & {[}τ{]} & {[}p{]} & Yes / No \\
Low-DTR extremes & {[}τ{]} & {[}p{]} & Yes / No \\
\bottomrule\noalign{}
\end{longtable}

\textbf{Interpretation:} Distribution tails {[}show independent trends /
follow mean signal{]}
\end{frame}

\section{Sensitivity: Seasonal vs.~Global
Thresholds}\label{sensitivity-seasonal-vs.-global-thresholds}

\begin{frame}{Sensitivity: Seasonal vs.~Global Thresholds}
\textbf{Do seasonal thresholds (per-month 90th/10th) alter conclusions?}

\begin{longtable}[]{@{}
  >{\raggedright\arraybackslash}p{(\linewidth - 6\tabcolsep) * \real{0.2500}}
  >{\raggedright\arraybackslash}p{(\linewidth - 6\tabcolsep) * \real{0.2500}}
  >{\raggedright\arraybackslash}p{(\linewidth - 6\tabcolsep) * \real{0.2500}}
  >{\raggedright\arraybackslash}p{(\linewidth - 6\tabcolsep) * \real{0.2500}}@{}}
\toprule\noalign{}
\begin{minipage}[b]{\linewidth}\raggedright
Metric
\end{minipage} & \begin{minipage}[b]{\linewidth}\raggedright
Global Threshold
\end{minipage} & \begin{minipage}[b]{\linewidth}\raggedright
Seasonal Threshold
\end{minipage} & \begin{minipage}[b]{\linewidth}\raggedright
Conclusion Unchanged?
\end{minipage} \\
\midrule\noalign{}
\endhead
Annual mean DTR & {[}Sig?{]} & {[}Sig?{]} & Yes / No \\
High extremes & {[}Sig?{]} & {[}Sig?{]} & Yes / No \\
Low extremes & {[}Sig?{]} & {[}Sig?{]} & Yes / No \\
\bottomrule\noalign{}
\end{longtable}

\textbf{Key finding:} Seasonal thresholds {[}do / do not{]} change
results, suggesting extremes are {[}defined by climatological context /
driven by absolute magnitude{]}.
\end{frame}

\section{Conclusion}\label{conclusion}

\begin{frame}{Conclusion}
\textbf{Hypothesis vs.~Finding} - \textbf{Expected (warming theory):}
DTR declining - \textbf{Found:} {[}YOUR RESULT{]} -
\textbf{Consistency:} {[}Consistent / Inconsistent / Inconclusive{]}

\textbf{Key Takeaways} 1. Mean DTR in St.~John's exhibits {[}no
significant trend / significant trend{]} 2. Extreme day frequencies show
{[}coupled / independent{]} trends 3. Maritime influence and natural
variability dominate over climate signal

\textbf{Limitations} - Regional averaging masks station-level
microclimatic differences - Short time series (26 years) limits power to
detect weak trends - Linear/monotonic trend assumptions; autocorrelation
not tested
\end{frame}

\begin{frame}
\end{frame}




\end{document}
